\hojaejercicio{Integración Compleja}
% Ejercicio 1
\ejercicio{}{
    Mediante traslación y giro del plano complejo y ponemos la intersección de ambas
    circunferencias en el origen y la otra intersección en el eje real positivo. 
    Luego, cogemos la transformación de Möbius:
    \[T(z) = \frac{1}{z}\]
    Entonces:
    \[T(0)=\infty, \quad T(R)=\frac{1}{R}, \quad T(r)=\frac{1}{r}\]
    Transforma nuestro conjunto en la banda $B=\{z \in \mathbb{C} : 1/R \leq \Re(z) < 1/r\}$.
    Entonces sea $T^1$ la inversa de T.\\
    Tenemos entonces que $T^1(B)$ es nuestro conjunto. Al ser la banda B convexo y $T^1$ un homeomorfismo,
    $T^1(B)$ es un conjunto simplemente conexo.
}

% Ejercicio 2
\ejercicio{Determina las siguientes integrales por la definición}{\begin{itemize}{
    \item[\textbf{a)}] $\int_{S(0;2)} \frac{z}{|z|} dz$:
    Parametrizamos $z = 2e^{it}, t \in [0, 2\pi]$. Entonces $|z|=2$ y $dz = 2ie^{it}dt$.
    \[ \int_0^{2\pi} \frac{2e^{it}}{2} (2ie^{it}) dt = 2i \int_0^{2\pi} e^{i2t} dt = 2i \left[ \frac{e^{i2t}}{2i} \right]_0^{2\pi} = 0 \]

    \item[\textbf{b)}] $\int_{S(0;1)} \bar{z} dz$:
    Parametrizamos $z = e^{it}, t \in [0, 2\pi]$. Entonces $\bar{z} = e^{-it}$ y $dz = ie^{it}dt$.
    \[ \int_0^{2\pi} e^{-it} (ie^{it}) dt = i \int_0^{2\pi} dt = i [t]_0^{2\pi} = 2\pi i \]

    \item[\textbf{c)}] $\int_{S(z_0;r)} (z-z_0)^n dz, n \in \mathbb{N}$:
    Parametrizamos $z = z_0 + re^{it}, t \in [0, 2\pi]$. Entonces $dz = ire^{it}dt$.
    \[ \int_0^{2\pi} (re^{it})^n (ire^{it}) dt = i r^{n+1} \int_0^{2\pi} e^{i(n+1)t} dt = i r^{n+1} \left[ \frac{e^{i(n+1)t}}{i(n+1)} \right]_0^{2\pi} = 0 \]

    \item[\textbf{d)}] $\int_{\gamma} z^2 dz$ (Camino cerrado):
    Dividimos en segmento real $\gamma_1$ y arco $\gamma_2$:
    \[ I_1 = \int_{-R}^R x^2 dx = \frac{2R^3}{3} \]
    \[ I_2 = \int_0^\pi (Re^{it})^2 iRe^{it} dt = iR^3 \left[ \frac{e^{i3t}}{3i} \right]_0^\pi = \frac{R^3}{3}(-1-1) = -\frac{2R^3}{3} \]
    Sumando ambas partes: $\int_{\gamma} z^2 dz = I_1 + I_2 = 0$.
    }\end{itemize}}

% Ejercicio 3
\ejercicio{}{\begin{enumerate}
    \item[\textbf{a)}] $\int_{\gamma} z^3 dz$:
    Como $f(z)=z^3$ es analítica, la integral depende solo de los extremos:
    \[ \int_0^1 z^3 dz = \left[ \frac{z^4}{4} \right]_0^1 = \frac{1}{4} \]

    \item[\textbf{b)}] $\int_{[-1+i, 1+i]} \operatorname{sen} z dz$:
    \[ [-\cos z]_{-1+i}^{1+i} = \cos(1-i) - \cos(1+i) = 2i \operatorname{sen}(1) \operatorname{senh}(1) \]

    \item[\textbf{c)}] $\int_{[-1+i, 1+i]} |z|^4 dz$:
    Parametrizamos $z(t) = t+i, t \in [-1, 1], dz=dt$:
    \[ \int_{-1}^1 (t^2+1)^2 dt = \int_{-1}^1 (t^4 + 2t^2 + 1) dt = \left[ \frac{t^5}{5} + \frac{2t^3}{3} + t \right]_{-1}^1 = \frac{56}{15} \]

    \item[\textbf{d)}] $\int_{S(0;1)} (\operatorname{Re} z)^2 dz$:
    Sea $z = e^{it}, dz = i e^{it} dt$:
    \[ \int_0^{2\pi} \cos^2 t (i e^{it}) dt = i \int_0^{2\pi} \cos^3 t dt - \int_0^{2\pi} \cos^2 t \operatorname{sen} t dt = 0 \]
    \end{enumerate}
    }

% Ejercicio 4
\ejercicio{Utiliza el criterio $M$ de Weierstrass para probar que si $(f_n)_{n=0}^\infty$ es una sucesión de funciones continuas en $\gamma^*$ tales que $|f_n(z)| \leq M_n$ para todo $z \in \gamma^*$, con $\sum_{n=0}^\infty M_n < \infty$, entonces:
\[ \int_\gamma f(z) dz = \sum_{n=0}^\infty \int_\gamma f_n(z) dz, \quad \text{donde } f(z) = \sum_{n=0}^\infty f_n(z) \]}{
    \begin{enumerate}
    \item \textbf{Convergencia Uniforme:} Dado que $|f_n(z)| \leq M_n$ para todo $z \in \gamma^*$ y la serie numérica $\sum M_n$ converge, el \textbf{Criterio M de Weierstrass} garantiza que la serie de funciones $f(z) = \sum_{n=0}^\infty f_n(z)$ converge absoluta y \textbf{uniformemente} en $\gamma^*$.
    
    \item \textbf{Continuidad:} Como cada $f_n$ es una función continua en el soporte de la curva $\gamma^*$ y la convergencia es uniforme, la función suma $f(z)$ es también una función continua en $\gamma^*$. Esto asegura que la integral $\int_\gamma f(z) dz$ está bien definida.
    
    \item \textbf{Acotación de la diferencia:} Definimos la suma parcial de la serie como $S_N(z) = \sum_{n=0}^N f_n(z)$. Para demostrar la igualdad, analizamos el límite de la diferencia entre la integral de la función límite y la integral de las sumas parciales:
    \[ \left| \int_\gamma f(z) dz - \sum_{n=0}^N \int_\gamma f_n(z) dz \right| = \left| \int_\gamma \left( f(z) - S_N(z) \right) dz \right| \]
    
    Aplicando la propiedad de acotación de integrales de línea ($|\int_\gamma g| \leq \max |g| \cdot L(\gamma)$):
    \[ \left| \int_\gamma (f(z) - S_N(z)) dz \right| \leq \max_{z \in \gamma^*} |f(z) - S_N(z)| \cdot L(\gamma) \]
    Donde $L(\gamma)$ representa la longitud (finita) del camino $\gamma$.

    \item \textbf{Compacidad y Límite:} Puesto que el camino $\gamma: [a,b] \to \mathbb{C}$ es continuo, su imagen $\gamma^*$ es un conjunto \textbf{compacto} en $\mathbb{C}$. Por el Teorema de Weierstrass (valores extremos), cualquier función continua en un compacto alcanza su máximo, por lo tanto:
    \[ \max_{z \in \gamma^*} |f(z) - S_N(z)| = \sup_{z \in \gamma^*} |f(z) - S_N(z)| \rightarrow 0 \]
    
    
    el límite de la diferencia es nulo. Por consiguiente, se cumple la propiedad de intercambio:
    \[ \int_\gamma f(z) dz = \sum_{n=0}^\infty \int_\gamma f_n(z) dz \]
\end{enumerate}
}

% Ejercicio 5
\ejercicio{
    Calcula  $\int_{\gamma} z^5 sen z dz$
}{
 Observamos primero que la función $f(z) = z^5 sen z$ es una \textbf{función entera} (analítica en todo $\mathbb{C}$), ya que es el producto de dos funciones enteras.

Analizamos el camino $\gamma$, compuesto por:


Dado que el camino empieza en $1-i$, pasa por $1+i$ y regresa a $1-i$, $\gamma$ es un \textbf{contorno cerrado}. 

Por el \textbf{Teorema de Cauchy-Goursat}, la integral de una función analítica sobre un contorno cerrado es cero:
\[ \int_{\gamma} z^5 \operatorname{sen} z dz = 0 \]
}

% Ejercicio 6
\ejercicio{}{}

% Ejercicio 7
\ejercicio{}{
    Tenemos:
\[
\int_{\gamma_1} \frac{1}{1+z^2} \, dz 
\quad \text{con } \gamma_1 = S(i,1)
\]

Sabemos que 
\[
\frac{1}{z+i} \in H(\mathbb{C} \setminus \{-i\}),
\]
que es un abierto conexo, por lo tanto, como 
$D(i,1) \subset \mathbb{C} \setminus \{-i\}$, 
tenemos:

\[
\int_{\gamma_1} \frac{1}{(z-i)(z+i)} \, dz 
= \int_{\gamma_1} \frac{g(w)}{w-i} \, dw
\]

Entonces, por la fórmula integral de Cauchy:
\[
\int_{\gamma_1} \frac{1}{(z-i)(z+i)} \, dz
= 2\pi i \, g(i) 
= 2\pi i \cdot \frac{1}{i+i} 
= \pi
\]
\[
\boxed{\int_{\gamma_1} \frac{1}{1+z^2} \, dz = \pi}
\]

\vspace{10pt}
Para el segundo vemos que la curva 
$\gamma_2 = \overline{IVU} \cup \overline{UVI}$ 
es homótopa a 0 en $\mathbb{C} \setminus \{i\}$, 
que es un abierto. \\[6pt]

Entonces, por el teorema de Cauchy para curvas homótopas:
\[
\int_{\gamma} f(z) \, dz = 0 
= \int_{\gamma_2} f(z) \, dz + \int_{\gamma_1} f(z) \, dz
\]
\[
\Rightarrow \int_{\gamma_2} f(z) \, dz 
= - \int_{\gamma_1} f(z) \, dz = -\pi
\]

\vspace{10pt}
Para el tercero ocurre algo similar que para el segundo:
\[
\int_{\gamma_3} f(z) \, dz = \pi
\]
}

% Ejercicio 8
\ejercicio{}{

\section*{8)}

\[
\frac{1}{2\pi i}\int_{S(0,1)} \frac{dw}{w^2(w-z)} \, dw \quad \text{para } |z|\neq 1.
\]

Si $|z|>1$ Entonces $\Rightarrow g(w)=\frac{1}{w^2(w-z)}\in H(\mathbb{C}/\{0,z\})$  
al ser cociente de holomorfas cuyo denominador no se anula en $\overline{D}(0,1)$.

Llamemos $h(w)=\frac{1}{w-z} \in H(\mathbb{C}/\{z\})$ y como $\overline{D}(0,1)\subset\mathbb{C}/\{z\}$  
entonces, por la fórmula integral de Cauchy para las derivadas:

\[
2\pi i \cdot 1 \cdot h'(0) =
\int_{\partial D(0,1)} \frac{h(w)}{(w-0)^2}\, dw.
\]

\[
 2\pi i \cdot \frac{1}{z^2}
\Rightarrow \frac{1}{2\pi i}\int_{S(0,1)} \frac{dw}{w^2(w-z)}\, dw
= \frac{1}{z^2}.
\]

Si $0<|z|<1$ entonces:  
vemos que $\tilde{\gamma}$ es homótopa a $0$ en $\mathbb{C}/\{0,z\}$ y por lo tanto

\[
\int f(w)\, dw = 0.
\]

\[
\Rightarrow \int_{\gamma} f(w)\, dw
= \int_{\gamma_1} f(w)\, dw
+ \int_{\gamma_2} f(w)\, dw.
\]

El procedimiento es igual que en los análogos ejercicios.

\begin{center}
\begin{tikzpicture}[scale=1.6, line cap=round, line join=round]

% Radios
\def\R{2}
\def\rA{0.8}
\def\rB{0.55}

% Punto z
\coordinate (Z) at (1.15,-0.35);

% Ejes
\draw[->] (-2.3,0)--(2.3,0);
\draw[->] (0,-2.3)--(0,2.3);

% Círculo grande r
\draw[thick] (0,0) circle (\R);
% flechas en el grande
\draw[->,thick] (0,0) ++(0:\R) arc (0:40:\R);
\draw[->,thick] (0,0) ++(90:\R) arc (90:130:\R);
\draw[->,thick] (0,0) ++(200:\R) arc (200:240:\R);
\draw[->,thick] (0,0) ++(300:\R) arc (300:340:\R);
\node at (1.6,1.3) {$r$};

% Círculo pequeño centrado r1
\draw[thick] (0,0) circle (\rA);
% flechas círculo centrado (CW)
\draw[<-,thick] (0,0) ++(0:\rA) arc (0:70:\rA);
\draw[<-,thick] (0,0) ++(200:\rA) arc (200:260:\rA);
\node at (0.9,-0.1) {$r_1$};

% Círculo en z r2
\draw[thick] (Z) circle (\rB);
% flechas círculo z (CW)
\draw[<-,thick] (Z) ++(20:\rB) arc (20:140:\rB);
\draw[<-,thick] (Z) ++(220:\rB) arc (220:300:\rB);
\node at (1.55,-0.85) {$r_2$};

% punto central y punto z
\fill (0,0) circle (1pt);
\fill (Z) circle (1pt);
\node[below right] at (Z) {$z$};

% Conexiones entre curvas (rectas con flechas en ambos sentidos)

% r ↔ r1
\draw[->,thick] (0,0)++(0:\rA) -- (0,0)++(0:\R);
\draw[->,thick] (0,0)++(0:\R) -- (0,0)++(0:\rA);

% r1 ↔ r2
\draw[->,thick] (0,-\rA) -- ($(Z)+(0,\rB)$);
\draw[->,thick] ($(Z)+(0,\rB)$) -- (0,-\rA);

% r ↔ r2
\draw[->,thick] (\R,0.2) -- ($(Z)+(\rB,0.1)$);
\draw[->,thick] ($(Z)+(\rB,0.1)$) -- (\R,0.2);

\end{tikzpicture}
\end{center}






}